\readerheading{Why the ratio conditions matter.}
The identity in the preceding proof combines the supported radial point
with one boundary point to produce a point on $r_0$. Since $a\le\varepsilon$,
its distance from $O$ is $\varepsilon/s\ge1/2$. A candidate also
contains $O$, so it must contain $M_0$. The other boundary point then demands
more of the candidate's left boundary trace than its signed-center geometry
allows. This is why the application needs a ratio estimate: it converts the
local supplier geometry into this same four-point contradiction.
